% c_cgs.tex — Exhaustive Simplified Guide to Writing .cgs Files
\chapter{Writing \.cgs Files: The Complete Authoring Guide}
\label{chap:cgs-authoring}

This section is the definitive reference for authoring \texttt{.cgs} (ComplexGitSync
specification) files. It covers both the simplified shorthand syntax suitable for
most use cases and the full advanced syntax for complex topologies, with practical
guidelines, complete examples, and common pitfalls.

% =======================================================================
% WHAT IS A .cgs FILE
% =======================================================================
\subsection{What is a \.cgs File?}

A \texttt{.cgs} file is a \textbf{TOML-formatted} specification that describes a tree
of Git repositories to be managed as a single logical unit. It is the only
hand-authored input in ComplexGitSync; all other documents (\texttt{.gts},
\texttt{.lgr}) are generated automatically at runtime.

\textbf{Key properties:}

\begin{itemize}
  \item \textbf{Offline and deterministic:} Parsing, normalization, and validation
        require no network access. The same file always produces the same
        internal representation.
  \item \textbf{Never modified at runtime:} \texttt{.cgs} files are read-only during
        execution. Snapshots (\texttt{.gts}) capture workspace state; they do not
        update the specification.
  \item \textbf{Two equivalent forms:} A concise shorthand syntax for simple cases,
        and a verbose inline-table syntax for advanced overrides. Both normalize
        to the same canonical internal document.
  \item \textbf{Validation is strict:} Unknown fields, malformed identifiers, duplicates,
        and missing required keys are rejected before any Git operations begin.
\end{itemize}

% =======================================================================
% MINIMAL .cgs: THE SIMPLIFIED FORM
% =======================================================================
\subsection{Minimal \.cgs: The Simplified Form}
\label{subsec:cgs-simplified}

For the vast majority of projects, a \texttt{.cgs} file needs only two top-level
keys: \texttt{project} (a string) and \texttt{repos} (a list of repository identifiers).

\subsubsection{Basic Structure}

\begin{lstlisting}
project = "my-project"

repos = [
    "github:owner/repo1",
    "gitlab:group/repo2",
    "codeberg:user/repo3",
]
\end{lstlisting}

\subsubsection{Repository Identifier Syntax}

Each repository identifier follows the pattern:

\begin{center}
\texttt{provider:owner/repository}
\end{center}

\begin{description}
  \item[provider] One of the supported Git providers: \texttt{github},
        \texttt{gitlab}, \texttt{codeberg}, or \texttt{custom}.
  \item[owner] The owner (user or organization/group) on the provider.
  \item[repository] The repository name itself.
\end{description}

Nested namespace paths (e.g., GitLab subgroups) are supported. The final path
segment is the repository name; all earlier segments join into
\texttt{project\_owner\_name}:

\begin{lstlisting}
repos = [
    "gitlab:my-group/my-subgroup/my-repo",  % 3-level namespace
    "github:my-org/my-repo",              % 2-level namespace
]
\end{lstlisting}

\textbf{Valid providers and their default hosts:}

\begin{center}
\begin{tabular}{lll}
\textbf{Provider} & \textbf{Host} & \textbf{Remotes Generated} \\ \hline
\texttt{github}   & \texttt{github.com}   & SSH and HTTPS \\
\texttt{gitlab}   & \texttt{gitlab.com}   & SSH and HTTPS \\
\texttt{codeberg} & \texttt{codeberg.org} & SSH and HTTPS \\
\texttt{custom}   & explicit URL required & SSH and HTTPS \\
\end{tabular}
\end{center}

% =======================================================================
% NORMALIZATION
% =======================================================================
\subsection{Normalization: What the Simplified Form Expands To}
\label{subsec:cgs-normalization}

The simplified shorthand is automatically \textbf{normalized} into a complete
canonical document. Every omitted field receives a deterministic default value.

\subsubsection{Default Values Applied During Normalization}

\begin{center}
\begin{tabular}{lll}
\textbf{Field} & \textbf{Default} & \textbf{Applied To} \\ \hline
\texttt{format\_version} & \texttt{"1.0"} & document table \\
\texttt{default\_branch} & \texttt{"main"} & project and each repo \\
\texttt{fallback\_branch} & \texttt{"main"} & each repo \\
\texttt{access\_protocol} & \texttt{"ssh"} & each repo \\
\texttt{nested\_config} & \texttt{"auto"} & each repo \\
\texttt{relative\_path} & repo name & each repo \\
\end{tabular}
\end{center}

\subsubsection{Normalization Example}

The following simplified input:

\begin{lstlisting}
project = "demo"

repos = [
    "github:acme/core",
    "gitlab:acme/utils",
]
\end{lstlisting}

Normalizes to this canonical form (conceptually):

\begin{lstlisting}
[document]
format_version = "1.0"

[project]
name = "demo"
default_branch = "main"

[[repos]]
gitprovider = "github"
project_owner_name = "acme"
project_name = "core"
repo_name = "core"
default_branch = "main"
fallback_branch = "main"
access_protocol = "ssh"
nested_config = "auto"
relative_path = "core"

[[repos]]
gitprovider = "gitlab"
project_owner_name = "acme"
project_name = "utils"
repo_name = "utils"
default_branch = "main"
fallback_branch = "main"
access_protocol = "ssh"
nested_config = "auto"
relative_path = "utils"
\end{lstlisting}

\textbf{Important:} The project's \texttt{default\_branch} propagates to all repos
unless explicitly overridden per-repo.

% =======================================================================
% INLINE TABLE SYNTAX: OVERRIDING DEFAULTS
% =======================================================================
\subsection{Inline Table Syntax: Overriding Defaults}
\label{subsec:cgs-inline}

When a repository needs non-default values for any field, use an \textbf{inline
table} instead of a plain string. Only the fields that differ from defaults
need to be specified.

\subsubsection{Common Override Scenarios}

\begin{lstlisting}
project = "demo"

repos = [
    % Non-default branch
    { repository = "github:acme/core", default_branch = "develop" },
    
    % Different default and fallback branches
    { repository = "gitlab:acme/libs", 
      default_branch = "next", 
      fallback_branch = "main" },
    
    % Custom placement path
    { repository = "github:acme/docs", 
      relative_path = "documentation" },
    
    % HTTPS access instead of SSH
    { repository = "codeberg:org/tool", 
      access_protocol = "https" },
    
    % Disable nested config discovery
    { repository = "github:acme/legacy", 
      nested_config = "disabled" },
    
    % Custom provider with explicit URL
    { repository = "custom:team/internal", 
      gitprovider_url = "https://git.internal.company.com" },
]
\end{lstlisting}

\subsubsection{Combined Overrides}

Multiple overrides can be combined in a single inline table:

\begin{lstlisting}
repos = [
    { repository = "github:acme/web", 
      relative_path = "apps/web", 
      default_branch = "production", 
      fallback_branch = "staging",
      access_protocol = "https",
      nested_config = "disabled" },
]
\end{lstlisting}

% =======================================================================
% ADVANCED TABLE SYNTAX: LEGACY FORM
% =======================================================================
\subsection{Advanced Table Syntax: Legacy Form}
\label{subsec:cgs-advanced}

The legacy advanced form uses explicit TOML tables. Fully supported but not
recommended for new files.

\begin{lstlisting}
[project]
name = "demo"
default_branch = "main"

[[repos]]
repository = "github:acme/core"

[[repos]]
repository = "gitlab:acme/utils"
default_branch = "develop"
relative_path = "lib/utils"
\end{lstlisting}

\textbf{When to use advanced tables:}

\begin{itemize}
  \item Project-wide settings in the \texttt{[project]} table
  \item When you prefer visual separation between entries
  \item For compatibility with old ComplexGitSync versions
\end{itemize}

\textbf{Note:} The \texttt{[project]} table can also be inline:
\texttt{project = \{ name = "demo", default\_branch = "main" \}}.

% =======================================================================
% REPOSITORY PLACEMENT AND PATH RESOLUTION
% =======================================================================
\subsection{Repository Placement and Path Resolution}
\label{subsec:cgs-placement}

\subsubsection{relative\_path: Where Each Repository Lives}

The \texttt{relative\_path} field determines where a repository is cloned relative
to its parent. Default: the repository name.

\begin{lstlisting}
% Default: repo cloned at <parent>/repo_name
repos = ["github:acme/core"]

% Explicit: repo cloned at <parent>/apps/core
repos = [
    { repository = "github:acme/core", relative_path = "apps/core" }
]
\end{lstlisting}

\subsubsection{Project Root Auto-Mount Rule}

If exactly one repository name matches the project name, its
\texttt{relative\_path} is automatically set to \texttt{"."}:

\begin{lstlisting}
project = "CGSil1"

repos = [
    "gitlab:CGS_test/CGSil1",  % Auto-mounted at "."
    "gitlab:CGS_test/CGSil2",  % At "./CGSil2"
    "github:flipoyo/CGSih1",   % At "./CGSih1"
]
\end{lstlisting}

Equivalent to:

\begin{lstlisting}
project = "CGSil1"
repos = [
    { repository = "gitlab:CGS_test/CGSil1", relative_path = "." },
    { repository = "gitlab:CGS_test/CGSil2", relative_path = "CGSil2" },
    { repository = "github:flipoyo/CGSih1", relative_path = "CGSih1" },
]
\end{lstlisting}

\subsubsection{Cross-Repository References}

A repository can reference another repository in the same tree using relative paths:

\begin{lstlisting}
% In CGSil2.cgs (nested under CGSil2/):
repos = [
    { repository = "github:flipoyo/CGSih1", 
      relative_path = "../CGSih1", 
      nested_config = "disabled" }
]
\end{lstlisting}

Here, \texttt{../CGSih1} resolves from \texttt{<CGSHOME>/CGSil2/} to
\texttt{<CGSHOME>/CGSih1}, referencing the sibling declared elsewhere.

% =======================================================================
% NESTED CONFIGURATION DISCOVERY
% =======================================================================
\subsection{Nested Configuration Discovery}
\label{subsec:cgs-nested}

The \texttt{nested\_config} field controls whether ComplexGitSync looks for
additional \texttt{.cgs} files inside each repository.

\begin{description}
  \item[\texttt{auto}] (default) Loads the sole root-level \texttt{*.cgs} file if
        one exists; multiple matches are rejected.
  \item[\texttt{disabled}] Does not inspect the repository for nested config.
  \item[a relative path] A path ending in \texttt{.cgs} loads that exact file.
\end{description}

\subsubsection{Example: Multi-Level Nested Topology}

\begin{lstlisting}
project = "parent-project"

repos = [
    "github:org/parent",
    { repository = "github:org/child-a", nested_config = "auto" },
    { repository = "github:org/child-b", nested_config = "disabled" },
    { repository = "github:org/child-c", nested_config = "config/prod.cgs" },
]
\end{lstlisting}

% =======================================================================
% ACCESS PROTOCOLS AND CUSTOM PROVIDERS
% =======================================================================
\subsection{Access Protocols and Custom Providers}
\label{subsec:cgs-protocols}

\subsubsection{SSH vs HTTPS}

Default: SSH. To use HTTPS:

\begin{lstlisting}
repos = [
    { repository = "github:owner/repo", access_protocol = "https" },
]
\end{lstlisting}

Generated remotes:

\begin{itemize}
  \item SSH: \texttt{git@github.com:owner/repo.git}
  \item HTTPS: \texttt{https://github.com/owner/repo.git}
\end{itemize}

\subsubsection{Custom Git Providers}

For providers not in the built-in list, use \texttt{custom} with explicit URL:

\begin{lstlisting}
repos = [
    { repository = "custom:team/project", 
      gitprovider_url = "https://git.internal.mycompany.com" },
]
\end{lstlisting}

\textbf{Important:} \texttt{gitprovider\_url} is \textbf{required} for \texttt{custom}.

\begin{lstlisting}
% INCORRECT - will fail validation
repos = ["custom:team/project"]

% CORRECT
repos = [
    { repository = "custom:team/project", 
      gitprovider_url = "https://git.internal.company.com" }
]
\end{lstlisting}

% =======================================================================
% BRANCHES
% =======================================================================
\subsection{Branches: default\_branch and fallback\_branch}
\label{subsec:cgs-branches}

\subsubsection{Branch Resolution Order}

When cloning or pulling:

\begin{enumerate}
  \item Repository's \texttt{default\_branch}
  \item Project's \texttt{default\_branch}
  \item Repository's \texttt{fallback\_branch}
  \item Project's \texttt{fallback\_branch} (defaults to \texttt{main})
\end{enumerate}

\subsubsection{Practical Branch Patterns}

\begin{lstlisting}
% All repos use "main" (default)
project = "my-project"
repos = ["github:owner/repo1", "github:owner/repo2"]

% All repos default to "develop"
project = { name = "my-project", default_branch = "develop" }
repos = ["github:owner/repo1", "github:owner/repo2"]

% Per-repo overrides
project = "my-project"
repos = [
    "github:owner/core",                    % Uses default "main"
    { repository = "github:owner/feature", default_branch = "dev" },
    { repository = "github:owner/stable", 
      default_branch = "release", fallback_branch = "main" },
]
\end{lstlisting}

% =======================================================================
% COMPLETE FIELD REFERENCE
% =======================================================================
\subsection{Complete Reference: All Recognized Fields}
\label{subsec:cgs-fields}

\subsubsection{Top-Level Fields}

\begin{center}
\begin{tabular}{llll}
\textbf{Field} & \textbf{Type} & \textbf{Required} & \textbf{Default} \\ \hline
\texttt{project} & string or table & yes & — \\
\texttt{repos} & list & yes & — \\
\texttt{document} & table & no & \texttt{\{format\_version = "1.0"\}} \\
\end{tabular}
\end{center}

\subsubsection{Project Table Fields}

\begin{center}
\begin{tabular}{llll}
\textbf{Field} & \textbf{Type} & \textbf{Required} & \textbf{Default} \\ \hline
\texttt{name} & string & yes & — \\
\texttt{default\_branch} & string & no & \texttt{"main"} \\
\end{tabular}
\end{center}

\subsubsection{Repository Fields}

\begin{center}
\begin{tabular}{llll}
\textbf{Field} & \textbf{Type} & \textbf{Required} & \textbf{Default} \\ \hline
\texttt{repository} & string & yes* & — \\
\texttt{gitprovider} & string & no & from \texttt{repository} \\
\texttt{project\_owner\_name} & string & no & from \texttt{repository} \\
\texttt{project\_name} & string & no & from \texttt{repository} \\
\texttt{repo\_name} & string & no & from \texttt{repository} \\
\texttt{default\_branch} & string & no & from project or \texttt{"main"} \\
\texttt{fallback\_branch} & string & no & \texttt{"main"} \\
\texttt{access\_protocol} & string & no & \texttt{"ssh"} \\
\texttt{relative\_path} & string & no & repo name (or "." if matches project) \\
\texttt{nested\_config} & string & no & \texttt{"auto"} \\
\texttt{gitprovider\_url} & string & no** & — \\ \hline
\end{tabular}
\end{center}

\begin{itemize}
  \item[*] Either \texttt{repository} (shorthand) or the combination of
        \texttt{gitprovider}, \texttt{project\_owner\_name}, \texttt{project\_name}
        is required.
  \item[**] \texttt{gitprovider\_url} is required when \texttt{gitprovider = "custom"}.
\end{itemize}

% =======================================================================
% PRACTICAL EXAMPLES
% =======================================================================
\subsection{Practical Examples}
\label{subsec:cgs-examples}

\subsubsection{Example 1: Simple Monorepo-Style Project}

\begin{lstlisting}
project = "my-app"

repos = [
    "github:myorg/my-app",    % Root (auto-mounted at ".")
    "github:myorg/core",     % At ./core
    "github:myorg/ui",       % At ./ui
    "github:myorg/api",      % At ./api
]
\end{lstlisting}

\subsubsection{Example 2: Mixed Providers with Custom Branches}

\begin{lstlisting}
project = "cross-provider"

repos = [
    "github:myorg/main-repo",
    { repository = "gitlab:mygroup/data", default_branch = "data-main" },
    { repository = "codeberg:user/tools", access_protocol = "https" },
]
\end{lstlisting}

\subsubsection{Example 3: Complex Topology with Nested Configs}

\begin{lstlisting}
project = "big-system"

repos = [
    "github:org/root",
    { repository = "github:org/frontend", 
      relative_path = "apps/frontend", nested_config = "auto" },
    { repository = "github:org/database", 
      relative_path = "infra/db", nested_config = "disabled" },
    { repository = "github:thirdparty/lib", 
      relative_path = "external/thirdparty-lib", nested_config = "disabled" },
]
\end{lstlisting}

\subsubsection{Example 4: Real-World (cawaqsviz)}

Based on \texttt{examples/cawaqsviz.cgs}:

\begin{lstlisting}
project = { name = "CaWaQS-Viz", default_branch = "main" }

repos = [
    { repository = "gitlab:cawaqs/gviz/cawaqsviz", 
      relative_path = ".", fallback_branch = "main" },
    { repository = "github:flipoyo/HydrologicalTwinAlphaSeries", 
      relative_path = "external/HydrologicalTwinAlphaSeries", 
      fallback_branch = "main", nested_config = "disabled" },
    { repository = "github:flipoyo/user_guide_CaWaQS-Viz", 
      relative_path = "docs/CWV_user_guide", 
      fallback_branch = "main", nested_config = "disabled" },
]
\end{lstlisting}

% =======================================================================
% CLI AUTHORING TOOLS
% =======================================================================
\subsection{Authoring with CLI Tools}
\label{subsec:cgs-cli-authoring}

You do not need to write \texttt{.cgs} files by hand.

\subsubsection{discover: Scan an Existing Checkout}

\begin{lstlisting}[language=bash]
$ pixi run cgitsync discover ~/work/myproject --write draft.cgs
$ pixi run cgitsync validate draft.cgs
\end{lstlisting}

\subsubsection{configure: Interactive Prompt}

\begin{lstlisting}[language=bash]
$ pixi run cgitsync configure --output myproject.cgs
\end{lstlisting}

\subsubsection{create-cgs: Direct CLI Authoring}

\begin{lstlisting}[language=bash]
$ pixi run cgitsync create-cgs \
    --project myproject \
    --repo github:owner/repo1 \
    --repo gitlab:group/repo2 \
    --output myproject.cgs
\end{lstlisting}

All three use the same offline pipeline as direct TOML parsing.

% =======================================================================
% VALIDATION AND COMMON ERRORS
% =======================================================================
\subsection{Validation and Common Errors}
\label{subsec:cgs-validation}

Always validate before use:

\begin{lstlisting}[language=bash]
$ pixi run cgitsync validate myproject.cgs
\end{lstlisting}

\subsubsection{Validation Checks}

\begin{itemize}
  \item Presence of required keys (\texttt{project}, \texttt{repos})
  \item Valid repository identifier syntax
  \item Known provider names or custom with URL
  \item \texttt{custom} providers have \texttt{gitprovider\_url}
  \item Valid \texttt{nested\_config} values
  \item Valid \texttt{access\_protocol} values
  \item No duplicate repository declarations
  \item Non-empty strings for all text fields
\end{itemize}

\subsubsection{Common Errors and Fixes}

\begin{description}
  \item[\texttt{project is required}] Add \texttt{project = "name"}.
  \item[\texttt{repos is required}] Add \texttt{repos = [...]}. 
  \item[\texttt{expected 'provider:owner/repository'}] Fix identifier syntax.
  \item[\texttt{gitprovider\_url is required for custom provider}]
        Add \texttt{gitprovider\_url = "https://..."}.
  \item[\texttt{duplicate repository}] Remove one or make \texttt{relative\_path}
        distinct.
  \item[\texttt{nested\_config must be auto, disabled, or a path}]
        Use valid value. Paths must end in \texttt{.cgs}.
\end{description}

% =======================================================================
% TEMPLATES AND STARTING POINTS
% =======================================================================
\subsection{Templates and Starting Points}
\label{subsec:cgs-templates}

\subsubsection{Template Files}

\begin{description}
  \item[\texttt{examples/template.cgs}] Copy for normal authoring.
  \item[\texttt{examples/normalized\_template.cgs}] Reference for canonical form.
\end{description}

\subsubsection{Quick Start Checklist}

\begin{enumerate}
  \item Copy \texttt{examples/template.cgs} to \texttt{<your-project>.cgs}
  \item Edit \texttt{project} name
  \item Add repositories to \texttt{repos} list
  \item Override defaults only when necessary
  \item Validate: \texttt{pixi run cgitsync validate <file>.cgs}
  \item Initialize: \texttt{pixi run cgitsync initialise <file>.cgs}
\end{enumerate}

% =======================================================================
% BEST PRACTICES
% =======================================================================
\subsection{Best Practices}
\label{subsec:cgs-best-practices}

\begin{itemize}
  \item \textbf{Start simple:} Use plain strings until you need overrides.
  \item \textbf{Use auto-mount:} Name one repo after project for automatic root.
  \item \textbf{Validate early and often:} After every edit.
  \item \textbf{Keep comments:} Document non-default values.
  \item \textbf{Avoid advanced tables:} Inline tables are more maintainable.
  \item \textbf{Use \texttt{nested\_config = "disabled"} judiciously:}
        Only when certain repo has no \texttt{.cgs}.
  \item \textbf{Prefer SSH:} Works well for most providers.
  \item \textbf{Be explicit with custom:} Always include \texttt{gitprovider\_url}.
\end{itemize}

% =======================================================================
% ROUND-TRIP GUARANTEE
% =======================================================================
\subsection{Internal Round-Trip Guarantee}
\label{subsec:cgs-roundtrip}

ComplexGitSync guarantees a semantic round-trip:

\begin{center}
\texttt{.cgs $\rightarrow$ CgsDocument $\rightarrow$ GitTree $\rightarrow$ CgsDocument $\rightarrow$ .cgs}
\end{center}

Parsing and normalizing a generated \texttt{.cgs} reproduces the canonical document.
Whitespace and comments need not match byte-for-byte, but semantics must.
